\begin{answers}{ch:agentic-rl}{Agentic RL：エージェントを強化学習で訓練する}
\item $c=0.3$ なら暗算、$c=0.1$ なら電卓が得です。
報酬の期待値は、式\eqref{eq:arl-reward}から、暗算が正答率そのままの $0.75$、電卓が正答率からコストを引いた値です。
$c=0.3$ では $0.98-0.3=0.68<0.75$、$c=0.1$ では $0.98-0.1=0.88>0.75$ です。
電卓で上がる正答率は $0.98-0.75=0.23$ なので、$c$ がこれより小さいときだけ、電卓が得になります。
\item アドバンテージは約 $(0.5,\,0.5,\,0.5,\,-1.5)$ です。
平均は $0.95$、平均との差は $0.05$ が 3 つと $-0.15$ です。
2 乗の和 $3\times0.05^2+0.15^2=0.03$ を $G-1=3$ で割ると $0.01$ で、標準偏差は $0.1$ です（分母の $10^{-4}$ まで入れると $0.4995$ と $-1.4985$）。
4 本とも正解で、3 本は暗算、1 本は電卓を 1 回使った軌跡です。
差は 0.2 しかないのに、アドバンテージは、1 本が不正解だった報酬 $(1,1,1,0)$ のときと同じ大きさになりました。
長所は、報酬の差が小さくても手がかりが消えず、報酬の目盛りが変わっても更新の大きさがそろうことです。
短所は、わずかな差も正誤の差と同じ強さで押してしまうことです。
全員が正解したグループでは、$c$ が $0.2$ でも $0.01$ でも、電卓を使った軌跡はほぼ同じ強さで確率を下げられます。
この例は、本文で引いた「標準偏差で割ると、易しすぎる問題と難しすぎる問題を強調しすぎる」という指摘\cite{liu2025understanding}の一つの形です。ほぼ全員が正解する易しいグループでは、標準偏差が小さいので、わずかな差が大きく拡大されます。
\item 軌跡ごとに、最初の役割トークン $a_1$ の確率の比 $w=\pitheta(a_1\mid s_1)/\mu(a_1\mid s_1)$ を、式\eqref{eq:arl-loss}の方策勾配の項に掛けます（式\eqref{eq:pg-importance}）。
$\mu$ は実際にサンプリングに使った分布で、\code{<|tool|>} と \code{<|assistant|>} では $\mu=(1-\varepsilon)\pitheta+\varepsilon/2$ です。
2 トークン目からはモデルが自分で選ぶので、比は 1 です。
「\code{65-46は？}」型で電卓を選んだ軌跡は、$\mu=0.8\times0.001+0.1\approx0.10$ なので $w\approx0.01$ になり、ほとんど数えられません。
偏りを取り除くと、平均では探索なしの更新に戻るので、探索で得た速さも失います。本章が許した偏りは、速さの代金です。
強制したトークンだけを学習から外すと、強制された電卓の軌跡が良い報酬を得ても、\code{<|tool|>} の確率は上がりません。
式の書き方は学べますが、「電卓を選ぶ」ことは、モデルが自分で選んだ軌跡からしか学べなくなります。
確率 0.001 の質問では、それはほとんど起きません。
\item 学習曲線は「RL のみ」に近いままで、コールドスタートの代わりにはなりませんでした（実測。乱数の種は 1 つ）。
学習前のモデルで、温度 1、$\varepsilon=0.2$ のロールアウトを 92,160 本作ると、電卓を使って正解した軌跡は 1,271 本（1.4\%）でした。
その 7 割は、\code{<|tool|>399<|end|>} のように、暗算した答えをそのまま電卓に渡した軌跡です。
質問どおりの式（\code{348+83} など）を書いた軌跡は、15 本しかありませんでした。
これに、コールドスタートと同じ数の掛け算と掛けて足す問題のお手本を足して SFT すると、加減算の電卓使用率は 42.8\% から 74.7\% に上がりましたが、正答率は 31.6\% が 32.3\% になっただけでした。
続けて $c=0$ で 73 ステップ強化学習をすると、学習中の電卓使用率は 7 割前後から 6 割台へ下がり、正答率は 2〜3 割台のままでした（\figref{fig:arl-curves}の実線は 8 割台と 6 割前後）。
貪欲の評価は正答率 44.7\%、電卓使用率 64.9\% で、正答率は RL のみの 43.5\% とほぼ同じです。
成功した軌跡を選んでも、学習前のモデルがほとんど試さない「正しい式を書く」振る舞いは集まらないからです（\figref{fig:arl-why-quit}）。
人が式を書いたお手本が教えていたのは、この部分でした。
\item 電卓を信用したまま、毎回 1 回だけ呼ぶことを学ぶと予想されます。
結果をそのまま答えれば 9 割は正解になり、暗算よりずっと得だからです。
誤りを疑って呼び直す、検算する、といった振る舞いは、お手本にも探索にも出てこないので、試されず、強められません。
実際に、電卓が確率 0.1 で 1〜9 だけずれた値を返す環境で、コールドスタート後のモデルから $c=0$ で 73 ステップ学習させました（乱数の種は 1 つ）。
加減算の電卓使用率は 100\%、呼び出しは 1 問に 1 回のままで、正しい電卓で測った正答率は 94.7\% でした（誤りのない環境で学習すると 92.8\%）。
評価でも電卓を 1 割誤らせると、正答率は 84.1\% に下がりました。誤った値を受け取った 110 問のうち 107 問で、その値をそのまま答えています。
同じ環境で $c=0.2$ を付けて学習しても、加減算の電卓使用率は 99.7\% でした。
強化学習が強めるのは、試した中で平均して得だった振る舞いです。信用するほうが平均して得なあいだは、信用は弱まりません。
疑うことを学ばせるには、確かめる手段と、それを試す機会を、環境やお手本の側で用意する必要があります。
\end{answers}
